%%═══════════════════════════════════════════════════════════════════
%% Lecture 01 — Introduction to Nuclear Physics
%% 77603 — Nuclear Physics
%% Date: [Not stated in transcript — Zoom session, first lecture of semester]
%% Lecturer: Dr. Moshe Friedman
%%═══════════════════════════════════════════════════════════════════

\section{Lecture 01 — Introduction to Nuclear Physics}
\label{sec:lec01}
\date{24 March 2026}


%%───────────────────────────────────────────────────────────────────
\subsection*{Overview}
%%───────────────────────────────────────────────────────────────────

Nuclear physics is the branch of physics that studies the atomic nucleus.
The nucleus is \emph{not} a point particle: it is a complex, quantum
many-body system composed of protons and neutrons (collectively called
\emph{nucleons}), bound together by three distinct interactions —
the electromagnetic force, the weak interaction, and the nuclear force.
Unlike classical many-body systems (e.g.\ gravitational $N$-body
simulations), there is no scale separation in the nucleus, making exact
solutions intractable even on the most powerful computers.
Compounding this difficulty, the nuclear force itself is not a
fundamental force of nature but rather a residual effect of the strong
(QCD) interaction acting between quarks inside nucleons, and no
closed-form expression for it exists.
This lecture introduces the scope and goals of the course, standard
nuclear notation, the Chart of Nuclides, and the characteristic units
used in nuclear physics.


%%───────────────────────────────────────────────────────────────────
\subsection*{Definitions}
%%───────────────────────────────────────────────────────────────────

\dfn{Nucleon}{A collective name for the proton and the neutron.
A free proton carries electric charge $+e$; a free neutron is
electrically neutral.
Both have nearly identical masses and are treated in the course as
two quantum states of a single object (anticipating isospin symmetry).}

\dfn{Atomic Nucleus}{A bound quantum many-body system consisting of $Z$ protons and $N$
neutrons (nucleons), with total nucleon number $A = Z + N$ (the
\emph{mass number}).
Standard notation:
\begin{equation}
  {}^{A}_{Z}\mathrm{X}
  \label{eq:lec01:notation}
\end{equation}
where $\mathrm{X}$ is the chemical symbol of the element.
Because the chemical symbol uniquely determines $Z$, the subscript
$Z$ is often omitted.
$N$ is almost never written explicitly in the notation, but is given
by $N = A - Z$.}

\dfn{Isotopes}{Two nuclei are \emph{isotopes} of the same element if they share the
same proton number $Z$ but have different neutron numbers $N$.
They occupy the same column in the nuclear chart (same $Z$, different $N$).}

\dfn{Isotones}{Two nuclei are \emph{isotones} if they share the same neutron number $N$
but have different proton numbers $Z$.}

\dfn{Isobars}{Two nuclei are \emph{isobars} if they share the same mass number
$A = Z + N$ while both $Z$ and $N$ individually differ.
Isobars lie on anti-diagonal lines in the Chart of Nuclides
(constant $Z+N$).}

\dfn{Nuclear Force / Nucleon--Nucleon Force}{The residual interaction acting between nucleons inside the nucleus.
It is \emph{not} a fundamental interaction; it arises as a residue of
the colour (QCD) force binding quarks inside each nucleon.
Because no first-principles derivation from QCD at nuclear
energy scales is currently tractable, the nuclear force is described
by \emph{effective models} (Yukawa-type potentials, chiral EFT, etc.).
Key empirical properties:
\begin{itemize}
  \item \textbf{Short range}: operates only within $r \lesssim 1$\,fm
    (order of the nucleon radius); can be modelled to zeroth order as
    a square well potential.
  \item \textbf{Strong}: for $r \lesssim 1$\,fm it dominates the
    Coulomb repulsion between protons.
  \item \textbf{Vanishes outside range}: to good approximation, the
    nuclear force is zero for $r \gg 1$\,fm.
\end{itemize}}

\dfn{Weak Interaction inside the Nucleus}{The weak interaction (preferably called the \emph{weak interaction}
rather than ``weak force'') is responsible for interconverting protons
and neutrons:
\begin{align}
  \beta^{-}\text{-decay:}&\quad n \to p + e^{-} + \bar{\nu}_{e}
  \label{eq:lec01:betaminus}\\
  \beta^{+}\text{-decay:}&\quad p \to n + e^{+} + \nu_{e}
  \label{eq:lec01:betaplus}
\end{align}
The weak interaction acts at a range shorter than the nuclear force
itself (it operates \emph{inside} the nucleon).}

\dfn{Chart of Nuclides (Nuclear Chart)}{A two-dimensional map of all known nuclei, with
\begin{itemize}
  \item $x$-axis: neutron number $N$,
  \item $y$-axis: proton number $Z$ (i.e.\ the element).
\end{itemize}
Each cell represents one nuclide ${}^{A}_{Z}\mathrm{X}$.
Colour coding (one common convention):
\begin{itemize}
  \item \textbf{Black}: stable nuclides (found naturally on Earth).
  \item \textbf{Pink/rose (right of black strip)}: $\beta^{-}$ emitters
    — too many neutrons; a neutron converts to a proton, moving one step
    diagonally toward stability.
  \item \textbf{Blue (left of black strip)}: $\beta^{+}$ emitters
    (or electron capture) — too many protons.
  \item \textbf{Yellow}: $\alpha$ emitters — nucleus ejects a
    ${}^{4}\mathrm{He}$ nucleus ($2p + 2n$); the daughter moves two
    steps down and two steps left in the chart.
  \item \textbf{Green}: fission — nucleus splits into two large fragments.
  \item \textbf{White/empty}: nuclides that either do not exist as bound
    systems or whose properties have not yet been measured.
\end{itemize}
The chart can be thought of as a topographic map where ``altitude''
represents energy; all decay modes correspond to motion toward the
black strip (the energy valley of stability).}

\dfn{Atomic Mass Unit ($u$ or $\mathrm{amu}$)}{The \emph{atomic mass unit} is defined so that the mass of
${}^{12}\mathrm{C}$ is exactly $12\,u$:
\begin{equation}
  1\,u \;\equiv\; \frac{1}{12}\,m({}^{12}\mathrm{C})
  \;\approx\; 931.5\;\mathrm{MeV}/c^{2}
  \;\approx\; 1\;\mathrm{GeV}/c^{2}.
  \label{eq:lec01:amu}
\end{equation}
In these units the masses of the free proton and neutron are
\begin{align}
  m_{p} &= 1.00728\;u, \label{eq:lec01:mp}\\
  m_{n} &= 1.00867\;u. \label{eq:lec01:mn}
\end{align}
The fractional mass difference $(m_{n}-m_{p})/m_{p} \approx 0.14\%$
underscores that proton and neutron are nearly identical objects
(two isospin states of the nucleon).
The electron mass is $m_{e} \approx 5\times10^{-4}\,u$, roughly
$1/2000$ of the nucleon mass, and is usually neglected in nuclear mass
budgets.}


%%───────────────────────────────────────────────────────────────────
\subsection*{Key Results}
%%───────────────────────────────────────────────────────────────────

\textbf{Nuclear size hierarchy.}  The characteristic length scales are

\begin{equation}
  r_{\text{quark}} \sim 0 \text{ (point-like)}
  \;\ll\;
  r_{\text{nucleon}} \sim 1\;\mathrm{fm}
  \;\ll\;
  r_{\text{nucleus}} \sim \text{few fm}
  \;\ll\;
  r_{\text{atom}} \sim 1\;\text{\AA} = 10^{-10}\;\mathrm{m}.
  \label{eq:lec01:scales}
\end{equation}
The nucleus is approximately $1/50{,}000$ the size of the atom.

\textbf{Nuclear binding energy scale.}
\begin{equation}\boxed{
  E_{\text{bind}} \sim \text{few}\;\mathrm{MeV \; per \; nucleon}
  \label{eq:lec01:bind_scale}
}\end{equation}
This is $10^{6}$ times larger than the energy scale of chemical bonds
($\sim\mathrm{eV}$), which explains both the enormous energy density
of nuclear reactions and the need for the dedicated units below.

\textbf{Stability line.}
For light nuclei the stable valley follows approximately
$N \approx Z$ (the $N=Z$ line).
For heavier nuclei the Coulomb repulsion between protons requires
$N > Z$ to maintain stability; beyond $N \approx 50$ the $N=Z$ line
lies entirely outside the stable region.


%%───────────────────────────────────────────────────────────────────
\subsection*{Derivations}
%%───────────────────────────────────────────────────────────────────

No quantitative derivations were performed in this introductory
lecture.  The lecturer indicated that derivations will begin in
Lecture 02 and beyond, after students have completed the assigned
reading (Krane, Chapter 3).


%%───────────────────────────────────────────────────────────────────
\subsection*{Examples}
%%───────────────────────────────────────────────────────────────────

\ex{Reading nuclear notation}{The nuclide ${}^{56}_{26}\mathrm{Fe}$ (iron-56) has:
\begin{itemize}
  \item $Z = 26$ protons (the subscript, or equivalently the chemical
    symbol Fe),
  \item $A = 56$ nucleons (the superscript),
  \item $N = A - Z = 30$ neutrons.
\end{itemize}
Its mass is approximately $56\,u$ (to within nuclear binding energy
corrections of order $\sim 0.1\%$ per nucleon).}

\ex{Identifying isotopes, isotones, and isobars}{Consider ${}^{12}\mathrm{C}$ and ${}^{13}\mathrm{C}$:
both have $Z=6$, but $N=6$ and $N=7$ respectively — they are
\emph{isotopes}.

Consider ${}^{12}\mathrm{C}$ ($Z=6, N=6$) and ${}^{11}\mathrm{B}$
($Z=5, N=6$): both have $N=6$ but different $Z$ — they are
\emph{isotones}.

Consider ${}^{14}\mathrm{C}$ ($Z=6, N=8$) and ${}^{14}\mathrm{N}$
($Z=7, N=7$): both have $A=14$ but differ in both $Z$ and $N$ —
they are \emph{isobars}.}

\ex{The NNDC interactive chart}{The National Nuclear Data Center (NNDC) hosts the interactive
Chart of Nuclides at \texttt{https://www.nndc.bnl.gov}.
Clicking on any nuclide reveals its half-life, decay modes, excited
energy levels, spin-parity $J^{\pi}$, binding energy, and much more.
The lecturer emphasised this as the primary numerical reference for
the entire course.}


%%───────────────────────────────────────────────────────────────────
\subsection*{Connections}
%%───────────────────────────────────────────────────────────────────

\begin{itemize}
  \item \textbf{Quantum Mechanics II (prerequisite)}: The lecturer
    explicitly stated that students should be familiar with the
    hydrogen atom, orbital structure, the harmonic oscillator, and
    finite square-well potentials from QM2.  These reappear in the
    shell model (Lecture $\sim$8).
  \item \textbf{Classical mechanics}: $N$-body gravitational systems
    (solar system) were cited as an analogy: even three-body classical
    systems cannot be solved analytically, let alone the quantum
    many-body nucleus.
  \item \textbf{Cosmological $N$-body simulations}: Mentioned as the
    classical many-body benchmark — tractable because of scale
    separation and a known force ($GMM/r^{2}$). The nucleus lacks both
    advantages.
  \item \textbf{QCD / Standard Model}: The nuclear force is a residue
    of the QCD colour force.  At laboratory nuclear energies QCD
    cannot be expanded perturbatively (the coupling is $O(1)$), so
    the nuclear force remains phenomenological.  High-energy nuclear
    experiments (e.g.\ heavy-ion collisions at the LHC) do probe
    perturbative QCD.
  \item \textbf{Ab-initio nuclear theory} (Nir Barnea group, HUJI):
    Mentioned explicitly; performs first-principles calculations
    starting from nucleon--nucleon potentials, currently limited to
    light nuclei ($A \lesssim 20$--30).
  \item \textbf{Nuclear astrophysics}: Stars are the natural environment
    where nuclear reactions occur at scale; the lecturer identified
    nuclear astrophysics as one of the major frontiers of contemporary
    nuclear physics.
  \item \textbf{Solid-state / condensed-matter physics}: The
    \emph{form factor} concept (topic of Lecture 02 reading) also
    appears in condensed-matter scattering; students should encounter
    it there later.
\end{itemize}


%%───────────────────────────────────────────────────────────────────
\subsection*{To Remember}
%%───────────────────────────────────────────────────────────────────

\begin{enumerate}
  \item The nucleus is a \textbf{quantum many-body system} — exact
    solutions are intractable; all nuclear models are effective models.
  \item Three interactions act inside the nucleus: \textbf{electromagnetic}
    (repulsive, Coulomb, long-range), \textbf{weak} (interconverts $p$
    and $n$, very short-range), and the \textbf{nuclear force}
    (attractive, short-range $\sim1\,\mathrm{fm}$, not derivable from
    first principles at nuclear energies).
  \item Standard notation: ${}^{A}_{Z}\mathrm{X}$, with
    $A = Z + N$.  Chemical symbol $\Leftrightarrow$ $Z$.
  \item The \textbf{unit of length} in nuclear physics is the
    \emph{femtometre} (fm, also called the ``fermi''):
    $1\,\mathrm{fm} = 10^{-15}\,\mathrm{m}$.
  \item The \textbf{unit of mass} is the \emph{atomic mass unit}:
    $1\,u \approx 931.5\,\mathrm{MeV}/c^{2}$, defined via
    $m({}^{12}\mathrm{C}) \equiv 12\,u$.
  \item Nuclear \textbf{binding energies} are of order
    $\sim\mathrm{MeV}$ per nucleon — $10^{6}$ times chemical bond
    energies.
  \item Isotopes share $Z$; isotones share $N$; isobars share $A$.
\end{enumerate}


%%───────────────────────────────────────────────────────────────────
\subsection*{Exam / HW Hints}
%%───────────────────────────────────────────────────────────────────

\begin{itemize}
  \item The exam will be \textbf{100\% written}, with no homework
    submission component.  However, Dr.\ Friedman stated that exam
    problems will be \emph{based on, and may include verbatim,
    problems from the homework sets} he distributes.  Solving all
    homework problems is therefore the primary exam preparation.
  \item Past exams from previous lecturers (e.g.\ Ami Leviatan) are of
    \textbf{limited use} because the course structure is substantially
    different under Dr.\ Friedman.
  \item Some homework problems will involve analysing \textbf{real
    nuclear data} (not only analytical calculations).  Familiarity
    with the NNDC website (\texttt{nndc.bnl.gov}) is expected.
  \item \textbf{Reading before Lecture 02 is mandatory}: Chapter 3 of
    Krane, \textit{Introductory Nuclear Physics}, with emphasis on
    \textbf{nuclear masses} and the \textbf{form factor} (first
    encounter with the concept).
\end{itemize}

%%═══════════════════════════════════════════════════════════════════
%% SUGGESTED MACROS (add to preamble.sty if not already present)
%%
%% \newcommand{\fm}{\,\mathrm{fm}}
%%     % femtometre: $r \sim 1\fm$
%%
%% \newcommand{\MeV}{\,\mathrm{MeV}}
%%     % megaelectronvolt: $E \sim 8\MeV$
%%
%% \newcommand{\amu}{u}
%%     % atomic mass unit symbol: $m_p = 1.00728\,\amu$
%%
%% \newcommand{\nuclide}[3]{{}^{#1}_{#2}\mathrm{#3}}
%%     % e.g. \nuclide{56}{26}{Fe} → ^{56}_{26}Fe
%%═══════════════════════════════════════════════════════════════════